\bookchapter{פרק 15: מדע הסוללה ואופטימיזציית חיי סוללה}

\chapterintro{סוללת טלפון (ליתיום-יון או ליתיום-פולימר) היא כמו ספוג אנרגיה: טעינה היא כמו שאיבת מים (אגירת חשמל), פריקה היא כמו סחיטת מים (צריכת חשמל). יש כמה מאפיינים שכדאי לזכור:}

\booksection{אחת, מה הסיפור של הסוללה? מחזורים ובריאות הסוללה}

: בניגוד לסוללות ישנות, אין לה אפקט זיכרון, לא חייבים לרוקן לפני הטעינה.

\begin{itemize}
  \item \textbf{פחד מחום}: טמפרטורה גבוהה היא הרוצח מספר אחת של סוללות.
  \item \textbf{מזדקנת}: אחרי שימוש ממושך, הקיבולת המרבית יורדת, זה תהליך טבעי.
  \item \textbf{מדדי ליבה: מספר מחזורים ובריאות הסוללה}מחזור אחד אינו טעינה אחת. למשל, היום צרכת 60\% וטענת למלא, מחר צרכת 40\% וטענת למלא, סך צריכה של 100\% נחשבת כמחזור אחד.
\end{itemize}

\booksubsection{בריאות הסוללה היא הקיבולת המרבית הנוכחית כאחוז מהקיבולת המקורית. מתחת ל-80\% תרגיש בבירור שחיי הסוללה מתקצרים.}

\begin{itemize}
  \item איך בודקים?
  \item iPhone: הגדרות > סוללה > בריאות סוללה וטעינה; אנדרואיד: בדרך כלל מחפשים "בריאות סוללה" בהגדרות.
  \item \textbf{שתיים, מיסטיקת טעינה? טעינה מהירה, טעינה אלחוטית וחוק הזהב}טעינה מהירה פוגעת בסוללה?
\end{itemize}

\booksection{שתיים, מיסטיקת טעינה? טעינה מהירה, טעינה אלחוטית וחוק הזהב}

\begin{itemize}
  \item \textbf{טעינה אלחוטית}: נוחה אך פחות יעילה, מתחממת יותר. מתאימה לטעינה לילית.
  \item \textbf{חוק הזהב לשימור}: נסה לשמור על הסוללה בין 20\% ל-80\%, אל תשאיר אותה רעבה (0\%) או דחוסה (100\%) כל הזמן.
  \item \textbf{שלוש, ניתוח מלא של פרוטוקולי טעינה מהירה: למה המטען שלך לא טוען מהר?}אולי נתקלת במצב הזה: הטלפון תומך בטעינה מהירה, אבל אחרי החלפת מטען הטעינה נהיית איטית מאוד. הסיבה היא שטעינה מהירה אינה רק "הספק גבוה", אלא דורשת שהמטען, הכבל והטלפון "ידברו באותה שפה" - אלה פרוטוקולי טעינה מהירה.
\end{itemize}

\begin{tcbinfo}{iPhone "80\% הגבלת טעינה" - חוק הזהב כהגדרה}
iPhone 15/16 (iOS 17+) - “20\%-80\%” כתכונה. עצירה ב-80\%, מאריך חיי סוללה. נתיב:

\texttt{הגדרות > סוללה > בריאות > טעינה > 80\%}

משמעות שלוש האפשרויות: (1)\textbf{ללא הגבלה}– טען עד 100\%; (2)\textbf{טעינת סוללה מותאמת}—לומד הרגלים, מאחר ל-100\% בסוף הלילה; 3)\textbf{מגבלת 80\%}—עוצר ב-80\%, הכי מגן. בבית/משרד: 80\%; בנסיעות: אופטימלי.

\end{tcbinfo}

\begin{tcbinfo}{מטען GaN – שותף מושלם לנסיעות עסקים}
שמת לב שמטענים קטנים? הגיבור:\textbf{GaN (גליום חנקן)}חומר חדש. יתרונות GaN על סיליקון:

\begin{itemize}
  \item \textbf{נפח קטן}:חצי או שליש מסיליקון, 65W קטן מ-5W של Apple.
  \item \textbf{יעילות גבוהה, חום נמוך}:הפסדי מתג נמוכים, 95\%+ יעילות, לא חם.
  \item \textbf{מתאים לנסיעות}:65W/100W רב-יציאות, טלפון+טאבלט+מחשב, קל.
  \item \textbf{המלצת רכישה}:PD GaN (65W+), Baseus, Anker, Nubia, Xiaomi. כבל E-Marker.
\end{itemize}

\end{tcbinfo}

\booksection{עקרונות בסיסיים של טעינה מהירה}

הספק טעינה = מתח × זרם (P = U × I). כדי להגדיל הספק, אפשר להעלות מתח, להעלות זרם, או את שניהם. פרוטוקולים שונים בחרו בנתיבים טכנולוגיים שונים.

\booksubsection{תהליך העבודה של טעינה מהירה:}

שלב לחיצת יד

: הטלפון והמטען מתקשרים דרך הכבל, מחליפים מידע על פרוטוקולים והספקים נתמכים.

\begin{enumerate}
  \item \textbf{שלב משא ומתן}: שני הצדדים קובעים את רמת ההספק המרבית המשותפת.
  \item \textbf{שלב טעינה}: המטען מוציא מתח וזרם לפי המוסכם, הטלפון מבצע ניהול חשמל.
  \item \textbf{שלב התאמה}: ככל שהסוללה מתמלאת, הטלפון מוריד בהדרגה את הספק הטעינה להגנת הסוללה.
  \item \textbf{2. פרוטוקולי טעינה מהירה מרכזיים בפירוט}USB PD (Power Delivery)
\end{enumerate}

\booksubsection{- מלך הפרוטוקולים האוניברסליים}

\textbf{מאפיינים}: תקן רשמי של ארגון USB-IF, בעל הכיסוי הרחב ביותר. PD 3.0 תומך עד 100W (20V/5A), PD 3.1 מורחב ל-240W (48V/5A).

\begin{itemize}
  \item \textbf{דרישות אביזרים}: 5V, 9V, 15V, 20V, PD 3.1 תומך גם ב-28V, 36V, 48V.
  \item \textbf{תאימות}: כל דגמי iPhone, Google Pixel, חלק מדגמי אנדרואיד הדגל, מחשבים ניידים, קונסולת Switch - כולם תומכים בפרוטוקול PD.
  \item \textbf{240W}: אוניברסליות גבוהה, מטען אחד מטעין מכשירים רבים; תומך ב-PPS (אספקת חשמל מתוכנתת), מאפשר ויסות מדויק של מתח וזרם.
  \item \textbf{חסרונות}: דורש ממשק Type-C ושבב ייעודי, עלות גבוהה יותר.
  \item \textbf{Qualcomm QC (Quick Charge)}- הוותיק של מחנה אנדרואיד
\end{itemize}

\textbf{QC 3.0}: תומך במתח מתכוונן 3.6V-20V בצעדי 0.2V, עד 18W.

\begin{itemize}
  \item \textbf{QC 4.0/4+}: תואם לפרוטוקול USB PD, תומך עד 27W, מוסיף טכנולוגיית טעינה דו-ערוצית.
  \item \textbf{QC 5.0}: תומך מעל 100W, 0-50\% תוך 5 דקות בלבד, שיפור משמעותי בחימום.
  \item \textbf{מתח גבוה מתכוונן}: רוב הטלפונים עם מעבדי Snapdragon תומכים, אך כל מותג בדרך כלל מעדיף את הפרוטוקול הפרטי שלו.
  \item \textbf{240W}- נציג הזרם הגבוה במתח נמוך
\end{itemize}

\textbf{נתיב טכנולוגי}: זרם גבוה במתח נמוך (החל מ-5V/4A), פחות חימום, יעילות גבוהה יותר.

\begin{itemize}
  \item \textbf{VOOC 1.0}: 20W (5V/4A), הושק ב-2014, פתח את עידן הטעינה המהירה.
  \item \textbf{SuperVOOC}: 50W-65W, משתמש בעיצוב סוללה כפולה בטור.
  \item \textbf{SuperVOOC 2.0}: 65W, ממלא סוללת 4000mAh תוך 30 דקות.
  \item \textbf{טעינת סופר פלאש 150W/240W}: טכנולוגיית הטעינה המהירה הסדרתית המהירה ביותר כיום.
  \item \textbf{מאפיינים}: דורש מטען וכבל ייעודיים (עם מגעים נוספים), לא תואם לפרוטוקולים אחרים.
  \item \textbf{דרישות אביזרים}- התפתחות הטעינה המהירה במתח גבוה
\end{itemize}

\textbf{FCP (Fast Charge Protocol)}: 9V/2A, 18W, תואם QC 2.0.

\begin{itemize}
  \item \textbf{SCP (Super Charge Protocol)}: תומך בהספק גבוה יותר, מ-22.5W דרך 66W ועד 100W.
  \item \textbf{מאפיינים טכניים}: משתמש בטכנולוגיית משאבת מטען, ממיר מתח גבוה לזרם גבוה במתח נמוך לסוללה, יעילות עד 98\%.
  \item \textbf{תאימות}: בלעדי לטלפוני Huawei/Honor, חלק מהמטענים תומכים בפרוטוקול PD.
  \item \textbf{240W}- עליית פרוטוקול מפיתוח עצמי
\end{itemize}

\textbf{שבב Surge P1}: שבב ניהול טעינה מהירה מפיתוח עצמי של Xiaomi, תומך בטעינה קווית של 120W.

\begin{itemize}
  \item \textbf{שבב Surge G1}: שבב ניהול סוללה, פועל יחד עם P1 לטעינה מהירה יעילה יותר.
  \item \textbf{מאפיינים טכניים}: משתמש בסוללה כפולה בטור + משאבת מטען, תומך בתאימות למספר פרוטוקולים.
  \item \textbf{תאימות}- שמרנית אך יציבה
\end{itemize}

\textbf{פרוטוקול PD}: תומך רק בפרוטוקול USB PD.

\begin{itemize}
  \item \textbf{נתיב טכנולוגי}: iPhone 8 ואילך תומכים בטעינה מהירה, עד כ-27W (iPhone 15 Pro Max יכול להגיע ל-27W+).
  \item \textbf{מאפיינים}: הספק שמרני יחסית, אך עקומת טעינה מאוזנת, בקרת חום מצוינת.
  \item \textbf{דרישות אביזרים}: דורש כבל USB-C to Lightning (או USB-C to USB-C) ומטען PD.
  \item \textbf{3. טבלת השוואת פרוטוקולי טעינה מהירה}פרוטוקול
\end{itemize}

\booksubsection{הספק מרבי}

\begin{center}
\begin{tabularx}{\linewidth}{|X|X|X|X|}
\hline
\textbf{נתיב טכנולוגי} & \textbf{תאימות} & \textbf{VOOC 1.0} & \textbf{240W} \tabularnewline
\hline
מתח גבוה/ריבוי רמות & בלעדי לסדרת OPPO & QC 5.0 & 100W+ \tabularnewline
מתח גבוה מתכוונן & מכשירי Qualcomm & VOOC/Dart & 240W \tabularnewline
זרם גבוה במתח נמוך & בלעדי לסדרת OPPO & SuperCharge & 100W \tabularnewline
משאבת מטען & Huawei/Honor & טעינה מהירה של Apple & \textasciitilde{}27W \tabularnewline
פרוטוקול PD & iPhone/iPad & 4. פרוטוקולי טעינה אלחוטית & לטעינה אלחוטית יש גם תקני פרוטוקול משלה: \tabularnewline
\hline
\end{tabularx}
\end{center}

\booksubsection{תקן Qi}

: נקבע על ידי WPC (ברית הטעינה האלחוטית), התקן האוניברסלי ביותר. Qi בסיסי תומך ב-5W-15W.

\begin{itemize}
  \item \textbf{תקן Qi2}: פורסם ב-2023, מציג מיקום מגנטי (דומה ל-Apple MagSafe), תומך ב-15W+.
  \item \textbf{פרוטוקולים פרטיים}: לכל מותג יש טעינה אלחוטית מהירה משלו, כמו טעינה אלחוטית 50W של Xiaomi, טעינה אלחוטית 50W של Huawei וכו', דורשים מטען ייעודי.
  \item \textbf{5. מדריך רכישת טעינה מהירה}בדוק התאמת פרוטוקול
\end{itemize}

\booksubsection{: ברכישת מטען, ודא שהוא תומך בפרוטוקול הטעינה המהירה של הטלפון שלך. מטענים של צד שלישי בדרך כלל מציינים את הפרוטוקולים הנתמכים.}

\begin{itemize}
  \item \textbf{בדוק את הכבל}: טעינה מהירה בהספק גבוה דורשת כבל ייעודי. כבלים רגילים עשויים לתמוך רק עד 60W.
  \item \textbf{בדוק אישורים}: משתמשי Apple מומלץ לבחור כבלים מאושרי MFi; משתמשי אנדרואיד לבחור מותג מקורי או צד שלישי מוכר.
  \item \textbf{בדוק חימום}: חימום הטלפון בזמן טעינה מהירה הוא נורמלי, אך אם הוא שורף את היד זה לא נורמלי, מומלץ להפסיק טעינה או להחליף מטען.
  \item \textbf{טיפים לבטיחות טעינה מהירה}הימנע ממשחק משחקים כבדים בזמן טעינה, יגרום לחימום מוגבר.
\end{itemize}

\begin{tcbwarn}{אל תשתמש במטענים וכבלים באיכות ירודה, עלולים ליצור סכנות בטיחות.}
\begin{itemize}
  \item הסר נרתיק טלפון כבד בזמן טעינה, כדי לעזור לפיזור חום.
  \item אם המטען או הטלפון מתחממים חריג או יש ריח מוזר, הפסק את השימוש מיד.
  \item ארבע, אלה הם טיפים אמיתיים לחיסכון בחשמל
  \item מסך (מלך צריכת החשמל)
\end{itemize}

\end{tcbwarn}

\booksection{ד. טכנולוגיית סוללה חדשה: סוללת אנודת סיליקון-פחמן}

ב-2024-2026, “סיליקון-פחמן” = סוללת דגל. שדרוג אנודה, פשוט:

\begin{itemize}
  \item \textbf{אנודת גרפיט מסורתית}:גרפיט, בוגר, 372mAh/g.
  \item \textbf{תוספת חומר סיליקון}:סיליקון 10x (4200mAh/g). סיליקון-פחמן = צפיפות גבוהה.
  \item \textbf{עליית צפיפות אנרגיה}:באותו נפח, קיבולת עולה \textasciitilde{}\textbf{15\%-20\%}. לכן דגלים 2024-2026 ב-7-9mm עם 5500-6000mAh.
\end{itemize}

\begin{tcbinfo}{אילו דגמים כבר משתמשים?}
\begin{itemize}
  \item \textbf{Honor}:Magic ראשון, Magic6 Pro/7/V3.
  \item \textbf{Huawei}:Xiaomi 14 Ultra/15.
  \item \textbf{vivo}:vivo X100 Ultra/X200 Pro.
  \item \textbf{Huawei/OPPO ואחרים}:2025-2026 דגלים, סיליקון-פחמן = סטנדרט.
\end{itemize}

\end{tcbinfo}

\textbf{השפעה מעשית על המשתמש}:היתרון הגדול ביותר של אנודת סיליקון-פחמן הוא “אותו נפח, קיבולת גדולה יותר” - הטלפון יכול להיות דק יותר ללא ויתור על חיי סוללה; או שהמכשיר נשאר באותו עובי, אך הסוללה עולה מ-4500mAh ל-5500mAh. במילים אחרות, אתה לא צריך לשלם “עובי” תמורת חיי סוללה. שים לב: סוללות סיליקון-פחמן דורשות עדיין אופטימיזציה של ניהול קור וטעינה מלאה עם שבבי טעינה מהירה, וזו אחת הסיבות שיצרנים ממליצים לשמור על 80\% טעינה לטווח ארוך.

\booksection{ה. אלו הן טיפים החיסכון בסוללה האמיתיים}

\begin{itemize}
  \item \textbf{אות חלש}: במעלית, במרתף, הטלפון יחפש אות בטירוף, יצרוך חשמל מהר.
  \item \textbf{5G}: אם האות לא טוב, עבור ידנית ל-4G, חיי הסוללה ישתפרו משמעותית.
  \item \textbf{שירותי מיקום}: לאפליקציות שאינן ניווט, הגדר הרשאת מיקום ל"רק בזמן השימוש".
  \item \textbf{הגבלת פעילות רקע}: iPhone בהגדרות > כללי > רענון אפליקציות ברקע; אנדרואיד בהגדרות סוללת אפליקציה.
  \item \textbf{בוא נבדוק: בדיקת רפואה לסוללה שלך}בדוק את הבריאות: מצא ורשום את הקיבולת המרבית של הסוללה שלך.
\end{itemize}

\begin{tcbpractice}{נתח את פרטי צריכת החשמל: הסתכל בתרשימים וברשימות האפליקציות בהגדרות הסוללה, מצא את האשמים בצריכת חשמל.}
\begin{enumerate}
  \item נסח אסטרטגיה: הגבל אפליקציות שצורכות חשמל חריג ברקע.
  \item פרק 16: מערכת חשבונות ושירותי ענן
  \item יש לי חבר, בשנה שעברה הטלפון שלו נגנב. הדבר הראשון שחשב עליו לא היה "כמה עולה הטלפון", אלא "האם אוכל לשחזר את התמונות שבפנים". לצערי, הוא לא הפעיל גיבוי ענן. הטלפון הוא חפץ חיצוני, אבל הנתונים בחשבון - אם הם אבדו, לא ניתן להחזירם. בפרק זה נדבר איך לרכוש ביטוח אמיתי לחיים הדיגיטליים שלך.
\end{enumerate}

\end{tcbpractice}
